\UseRawInputEncoding
\documentclass{article}

% if you need to pass options to natbib, use, e.g.:
%     \PassOptionsToPackage{numbers, compress}{natbib}
% before loading neurips_2026

% The authors should use one of these tracks.
% Before accepting by the NeurIPS conference, select one of the options below.
% "default" for submission
\PassOptionsToPackage{numbers,sort&compress}{natbib}
\usepackage{neurips_2026_vericode}

% After being accepted, the authors should add "final" behind the track to compile a camera-ready version.
% \usepackage[main, final]{neurips_2026_vericode}

% "preprint" option is used for arXiv or other preprint submissions
 % \usepackage[preprint]{neurips_2026_vericode}

% to avoid loading the natbib package, add option nonatbib:
%    \usepackage[nonatbib]{neurips_2026_vericode}

\usepackage[utf8]{inputenc} % allow utf-8 input
\usepackage[T1]{fontenc}    % use 8-bit T1 fonts
\usepackage{hyperref}       % hyperlinks
\usepackage{url}            % simple URL typesetting
\usepackage{booktabs}       % professional-quality tables
\usepackage{amsfonts}       % blackboard math symbols
\usepackage{nicefrac}       % compact symbols for 1/2, etc.
\usepackage{microtype}      % microtypography
\usepackage{xcolor}         % colors
\usepackage{listings}      % for code listings
\usepackage{graphicx}


\lstset{
    language=Haskell, %% closest to Lean I guess
    basicstyle=\ttfamily\small,
    breaklines=true,
    columns=fullflexible,
    keepspaces=true,
    showstringspaces=false
}

\newcommand{\code}{\lstinline}
\newcommand{\mathlib}{\code{Mathlib}}

\usepackage{xcolor}
\definecolor{keywordcolor}{rgb}{0.7, 0.1, 0.1}   % red
\definecolor{tacticcolor}{rgb}{0.0, 0.1, 0.6}    % blue
\definecolor{commentcolor}{rgb}{0.4, 0.4, 0.4}   % grey
\definecolor{symbolcolor}{rgb}{0.0, 0.1, 0.6}    % blue
\definecolor{sortcolor}{rgb}{0.1, 0.5, 0.1}      % green
\definecolor{attributecolor}{rgb}{0.7, 0.1, 0.1} % red
\definecolor{varcolor}{rgb}{0.42, 0.20, 0.56}    % muted violet, pseudocode variables
\definecolor{fncolor}{rgb}{0.05, 0.35, 0.55}     % steel blue, pseudocode procedures


\def\lstlanguagefiles{lstlang0.sty,lstlang1.sty,lstlang2.sty,lstlang3.sty,lstlean.tex}
\lstset{
    language=lean,
    inputencoding=utf8,
    extendedchars=true,
    escapeinside={{!}{!}},
    % Mapping the groups to your colours:
    keywordstyle=[1]\color{keywordcolor}\bfseries, % Keywords
    keywordstyle=[2]\color{sortcolor},             % Sorts (Type, Prop)
    keywordstyle=[3]\color{tacticcolor},           % Tactics (rw, simp)
    commentstyle=\color{commentcolor}\itshape,
    stringstyle=\color{keywordcolor},
    % Formatting tweaks
    sensitive=true,
    breaklines=true,
    showstringspaces=false
}

% \newcommand{\code}[1]{\lstinline[language=Haskell]{#1}}


\lstdefinestyle{pseudo}{
  language={},
  basicstyle=\ttfamily\small,
  morekeywords={for,each,in,if,return,and,or,not,of,every,union},
  keywordstyle=\color{keywordcolor}\bfseries,
  emph={usedVars,boundVars,nonlocals,comprehensionAndLambdaTargets,params},emphstyle={\color{fncolor}\bfseries},
  emph=[2]{free,caps,mutated,ordered,threaded,inner,outer},
  emphstyle=[2]{\color{varcolor}},
  comment=[l]{//},commentstyle=\color{commentcolor}\itshape,
  mathescape=true,breaklines=true,columns=fullflexible,keepspaces=true,
  numbers=left,numberstyle=\tiny\color{gray},numbersep=8pt,xleftmargin=14pt
}
\lstdefinestyle{lean}{language=lean,basicstyle=\ttfamily\small,breaklines=true,columns=fullflexible,keepspaces=true,showstringspaces=false,extendedchars=true}
\lstdefinestyle{py}{language=Python,basicstyle=\ttfamily\small,breaklines=true,columns=fullflexible,keepspaces=true,showstringspaces=false}

%% \newcommand{\code}{\texttt} % placeholder till we have proper Lean code support.
\newcommand{\todo}[1]{\textcolor{red}{TODO: #1}}
\newcommand{\suggest}[1]{\textcolor{blue}{#1}}
\newcommand{\leftsquigarrow}{\reflectbox{$\rightsquigarrow$}}

% Note. For the workshop paper template, both \title{} and \workshoptitle{} are required, with the former indicating the paper title shown in the title and the latter indicating the workshop title displayed in the footnote. 
\title{PastaLean: Deterministic transpilation from Python to Lean for Verification}


% The \author macro works with any number of authors. There are two commands
% used to separate the names and addresses of multiple authors: \And and \AND.
%
% Using \And between authors leaves it to LaTeX to determine where to break the
% lines. Using \AND forces a line break at that point. So, if LaTeX puts 3 of 4
% authors names on the first line, and the last on the second line, try using
% \AND instead of \And before the third author name.


\author{%
  Anirudh Gupta\\
  Indian Institute of Science\\
    Bengaluru \\
  \texttt{anirudhgupta@iisc.ac.in} \\
  % examples of more authors
  \And
  Petros Markopoulos \\
  University of California \\
  San Diego \\
  \texttt{pmarkopoulos@ucsd.edu} \\
  \AND
  Siddhartha Gadgil \\
  Indian Institute of Science \\
  Bengaluru \\
  \texttt{gadgil@iisc.ac.in} \\
}


\begin{document}


\maketitle


\begin{abstract}
  We present \emph{PastaLean}, an extensible system for \emph{deterministic} translation from Python to Lean via the Abstract Syntax Tree of Python. This allows for verification of Python code, especially auxiliary code used by AI. \suggest{Our tool runs provides coverage of > 92\% Leetcode, 100\% Humanvalplus and > 94\% LivecodeBench groundtruth python files and correctness of about 98\% test cases passing. }
\end{abstract}


\section{Introduction}

The rapid evolution of large language models (LLMs) has transformed them from systems primarily designed for natural-language understanding into increasingly capable programming agents that can generate, reason over, and manipulate executable code \cite{jiang2026codegeneration}. Agentic Artificial Intelligence (AI) systems make extensive use of Python for calculations, numerical experiments, data analysis, and other tasks. However, if the generated code is incorrect, the conclusions drawn from it are also incorrect. Formal verification promises something stronger, a machine-checked guarantee that a program meets a specification on \emph{all} inputs. Yet the languages in which such guarantees are routinely established such as Lean~4~\cite{moura2021lean4}, Coq, Dafny~\cite{leino2010dafny}, F$^\star$~\cite{swamy2016fstar}, Verus~\cite{lattuada2023verus} are not the languages practitioners actually use. The code that needs checking is usually written in Python.

To address this issue, we introduce \emph{PastaLean}, a system that provides \emph{deterministic} translation from Python to Lean. This enables verification of Python code.

PastaLean implements a recursive, extensible framework for transpiling Python code to Lean. Python code is first parsed into an Abstract Syntax Tree (AST) using Python's built-in \code{ast} module. The AST is then traversed recursively, and each node is translated into JSON using a set of translation functions with dynamic dispatch. The JSON is then converted into Lean code recursively using an extensible collection of translation functions based on Lean meta-programming. Each of these functions has output typed Syntax in Lean, so the compiler ensures that the output is syntactically correct. Further, a front end is used to verify that the output compiles. 

To verify the correctness of the translated Lean function, we use AI to generate a specification in Lean. The specification is translated independently, and we can have multiple special cases for additional verification. While the specification may be incorrect, it is highly unlikely that both the specification and the translation are incorrect in a manner that makes them consistent with each other.

To ensure correctness of the Python code, it is crucial that the translation is deterministic -- an LLM may translate incorrect Python code to correct Lean code. The verification depends on the correctness of the translation code, but this is a fixed, open-source code base that can be verified independently. We also use differential testing to enhance trust in the translation. Further, an error in translation is highly unlikely to translate incorrect Python code to correct Lean code, and our main goal is to ensure that auxiliary code used by LLMs is correct.

Agentic AI systems typically generate and run a lot of exploratory code, only a small fraction of which is used in the final output. Only the code that is used in the final output needs to be verified. This means the verification adds only a small overhead to the overall computation. Furthermore, verification can be done in parallel with the exploration, so it does not slow down the overall computation.

Our use of Python, rather than directly generated Lean code, is motivated by its familiarity to both developers and LLMs. An alternative approach is to use a system like Dafny to generate code directly in a verified language. However, in that case if a specification is mistranslated, the generated code may be correct with respect to the specification, but incorrect with respect to the intended behavior.

\todo{Can we write "Our main contributions include?" few pointers}

\subsection{Related Work}

The work closest to PastaLean is Aeneas, which translates safe Rust, via the LLBC intermediate representation, into pure functional models for verification in Lean and other proof assistants~\cite{aeneas}. PastaLean follows the same broad strategy of verifying source programs through functional translation. However, Python lacks the static type and ownership information exploited by Aeneas and requires support for dynamic typing, mutation, exceptions, and its object model. The rocq-of-python project also translates Python into a proof assistant, targeting Rocq and using a dedicated representation of Python semantics~\cite{rocq-of-python}. PastaLean instead produces verification-friendly shallow embeddings in Lean, allowing generated definitions to use Lean and Mathlib\cite{mathlib2020} libraries directly.

Other systems verify Python without translating it into an interactive theorem prover. Nagini translates statically typed Python into Viper and uses SMT-based automation for modular verification~\cite{nagini}, while CrossHair uses symbolic execution to search for counterexamples to Python contracts~\cite{crosshair}. Dafny takes the complementary approach of writing programs directly in a verification-aware language~\cite{dafny}. PastaLean is distinguished by starting from ordinary Python and producing Lean code through a deterministic, extensible translator.

Recent LLM-based systems further motivate this design. PyVeritas uses an LLM to translate Python to C before applying bounded model checking~\cite{pyveritas}, while VeriBench evaluates agents that generate Lean implementations, specifications, and proofs from Python~\cite{veribench}. In PastaLean, translation of the executed program is deliberately kept outside the LLM: AI may generate the specification and proof, but a fixed translator preserves the connection to the original Python code, and Lean checks the resulting proof.
\section{Transpilation and Verification}

Because of the nature of Python -- being dynamically typed, having mutable state, allowing for multiple returns, being object-oriented etc. -- there are many subtleties in transpiling Python to Lean. Further, there are choices to be made in transpiling Python to Lean, and different choices affect the ease of verification.

We outline our basic approach to transpilation, and then discuss how we address many of the issues.

\subsection{Recursive Transpilation Framework}

Given Python code, we first parse it into an Abstract Syntax Tree (AST) using Python's built-in \code{ast} module. The AST is then traversed recursively, and each node is translated into JSON using a set of translation functions with dynamic dispatch.

The core of our approach is to recursively build Lean code using meta-programming in Lean, using functions registered in an environment extension. Concretely, the Json corresponding to a Python expression, function or module is translated into Lean by calling the function \code{getCode (json: Json) (kind: SyntaxNodeKind) : PygenM <| TSyntax kind} which takes as input the Json and the kind of syntax node to be generated, and returns a monadic action that produces the corresponding Lean code. 

The Monad \code{PygenM} is an abbreviation (i.e., a reducible definition) for \code{StateT PyGen.State TermElabM}, i.e., a state monad transformer over the \code{TermElabM} monad, which is the monad used for elaborating terms in Lean. The state contains information relevant to the generated Lean code. For instance, the first occurrence of a mutable variable \code{n} has to be translated (within a \code{do} block) into a \code{let mut n := ...} declaration, while subsequent occurrences of the variable are translated into \code{n := ...}. The state keeps track of which variables have been declared as mutable, so that subsequent occurrences can be translated correctly.

The function \code{getCode} expects a string as the field "node\_type" in the Json, which is used to look up the corresponding translation function in the environment extension \code{pygenExt}. A functions corresponding to the node type "type" is added to the environment extension by annotation with the attribute \code{[pygen type]}.
These functions can call \code{getCode} recursively to translate sub-expressions.

\suggest{From the Json we also obtain a function signature/function call broken down. We reimplement the Python Standard API within Lean using its standard Libraries and Mathlib. The arguments are passed to these Lean functions along with any type annotations if given. More details in Appendix.} \todo{Add an appendix section for this}

\subsubsection{Transpilation targets}

We map Python into a \emph{shallow embedding} in Lean, i.e., we use Lean's own types and functions to represent Python types and functions wherever possible. This is to allow us to use Lean's existing libraries for verification. Indeed, we use two shallow embeddings of Python in Lean, one for verifiability and one for computability, so in general we get two translations that differ slightly. One of these generates computable Lean code that can be executed, for instance using \code{Float} (which is not completely specified and does not have many theorems that are proved, or even true), while the other generates Lean code that is more amenable to verification, using \code{Real} in place of \code{Float}, preferring \code{List} over \code{Array}, and using standard Lean functions as far as possible. The former is used to verify the correctness of the latter, and to run the code for testing.
\todo{Do we need to address the differences between the two and where they might matter? e.g. if you're doing numeric computations (which is not rare), the two might differ significantly, and, in fact, the difference might be exactly the part that would benefit from verification.}

Furthermore, we attempt to generate the most \emph{verification friendly} form of code possible. For instance, we prefer to use the functional style and recursion. Where we need to use a Monad only to support mutability, \code{while} loops, Exceptional Handling we use the \code{Id} Monad. We use custom state and exception monads rather than \code{IO} where possible since \code{IO} supports too many side-effects.

\subsection{Supporting Python}

We sketch below the main issues in transpiling Python to Lean, and how we address them.

\subsubsection{Arithmetic Operations}\label{sec:arith}

\suggest{
Python's arithmetic operators are polymorphic, meaning the same operator can work with different data types depending on the operands. For example, \code{+} performs numeric addition for integers and concatenation for strings, while Boolean values behave like integers (True + True = 2, False + 3 = 4). We model each as a typeclass whose \emph{result} type is an \code{outParam}: e.g. \code{PyHAdd α β (outParam γ)}(with a shorthand notation $+_p$ ) and analogously \code{PyHSub}($+_p$), \code{PyHMul}($*_p$), \code{PyHDiv}($/_p$), \code{PyModulo}($\%_p$), etc. \\
Operations follow the numeric widening hierarchy (Section \ref{numeric-widening}), with the result taking the wider operand's type. To model this in Lean, we create Typeclasses for each operation with instances created for different operand types to support internally the polymorphic behavior of Python Operators. We implement this widening with a single class \code{PyNumJoin} \ref{app:arith}, which records how any two numeric types convert into a common one, so that each operator needs only one generic instance rather than one per pair of types.}

\begin{lstlisting}[
caption={The \code{PyHAdd} class, its notation, and representative instances.},label={lst:pyhadd}, mathescape=true]

class PyHAdd (α β : Type) (γ : outParam Type) where
  hAdd : α → β → γ

infixl:65 " $+_p$ " => PyHAdd.hAdd

@[default_instance] instance {α β γ} [HAdd α β γ] : PyHAdd α β γ where hAdd := HAdd.hAdd
instance (priority := high) : PyHAdd Int Bool Int where hAdd a b := a + pyBoolToInt b
instance : PyHAdd String String String where hAdd := String.append
\end{lstlisting}

\suggest{
Figure \ref{lst:pyhadd} defines the \code{PyHAdd} type class, which establishes Python-style addition semantics. The default instance delegates to Lean's native \code{HAdd} for applicable heterogeneous operations(see Section \ref{app:arith} for details of why ordering of instance is important). It is important that these operators are correct with Lean's internal arithmetic which is highlighted more in Section \ref{arith-correct}. We also elaborate more on how a value such as \code{float('inf')} is implemented in Section \ref{floatinf}.}

\subsubsection{Variables and Mutability}
Variables in Python are mutable, meaning that their values can be changed after they are created. In Lean, variables are immutable by default, meaning that once a value is assigned to a variable, it cannot be changed. However, Lean does support mutable variables in monads through the use of the \code{let mut} construct, which allows us to create mutable variables.

In simple cases, mutability can be translated simply to reassignments in Lean. This can be done in functional code, which is thus easy to verify. We make a first-pass attempting to translate Python code into functional code in Lean. In this pass, we use reassignments in Lean to model mutability in Python.

If this fails, we use a monadic translation by generating a sequence of \emph{do elements} in Lean and wrap this is a \code{do} block -- in the simplest cases using just \code{Id.run do ...} to run the \code{do} block in the \code{Id} monad.

While Python has a uniform syntax for assignment and modification of variables, for a mutable variable \code{x} in Lean we need to use \code{let mut x :=} for the first assignment, and then \code{x :=} for subsequent modifications, both of which differ from the syntax for assignment/reassignment, \code{let x :=}. We keep track of which variables have been declared as mutable, so that subsequent occurrences can be translated correctly. \suggest{In case of Type mutation of a variable, we follow an SSA based approach of renaming the variables. See \ref{sec:ssa} for more details.}

\todo{We should cite "Do unchained"}

\subsubsection{Loops, State, Exceptions}\label{sec:loops}

When the Python code has \code{while} \suggest{and \code{for}} loops, or has state or throws exceptions, we use a monadic translation. We can build a monad supporting state and exceptions using the monad transformers \code{StateT} and \code{ExceptT}. We can then translate the Python code into a \code{do} block in Lean, using the monad we have built. Such a Monad is supporeted by Lean's \emph{mvcgen} and \emph{vcgen} tools, which can be used to verify the code. \suggest{We support Python's \code{try}, \code{except}, \code{finally} syntax for catching Exceptions follow a similar code style.}

Alternatively, we can use the \code{IO} monad, which supports state and exceptions, but also supports many other side-effects. This makes it harder to verify the code, since we need to reason about all the side-effects that the code may have. We prefer to use a custom monad that only supports state and exceptions, so that we can reason about the code more easily.

\subsubsection{Python's Standard functions}

\todo{talk about how we manually implement these}

\subsubsection{Polymorphic Functions}\label{sec:poly}

\suggest{A polymorphic function accepts arguments of several types, and Python has two different sources of it (Listing \ref{lst:polyeg}). Library functions such as \code{len}, \code{min} and \code{max} support a fixed, known set of types, and for these Lean's own type classes are exactly the right tool: each operation is a protocol class with one instance per supporting type (\code{PyLen} for \code{len}, \code{PyGetItem} for subscripting, \code{PyIterable} for iteration), the operand's type selects the instance during elaboration, nothing is boxed, and the result is as provable as any monomorphic definition. A user-written \code{add} declares no types at all, so there is no fixed set of types to write instances for and type inference has to decide instead.}

\begin{lstlisting}[language=python, caption={Two sources of polymorphism in Python.}, label={lst:polyeg}, breaklines=true]

# library functions: a fixed set of supported types
len("apple")          # 5
len([1, 2, 3])        # 3
min("a", "b")         # "a"
max(1, 2)             # 2

# a user-defined function with no declared types at all
def add(a, b): return a + b
add(1, 2)             # 3
add("pasta", "lean")  # "pastalean"
\end{lstlisting}

\suggest{Our Type inference engine (Section \ref{sec:typeinfer}) assigns such a parameter a concrete type wherever it can, and the generated definition is then ordinary monomorphic Lean. Where it genuinely cannot, the parameter falls back to \code{PyAny}, a tagged union that dispatches on the runtime tag; boxing costs provability, which Appendix \ref{app:pyany} discusses.}

\subsubsection{Classes and Python OOP} \label{sec:oop}

\suggest{ We translate a Python \code{Class} into a Lean \code{structure} paired with a \emph{namespace} of functions. Class attributes map directly to structure fields, the constructor becomes \code{ClassName.new}, and methods translate to standalone functions that explicitly take \code{self} as an argument. To handle state changes in a purely functional way, we use \emph{value semantics}: modifying \code{self} simply returns an updated object, which is then rebound at the call site. Finally, Python dunder methods (e.g., \code{__add__}) are mapped to Lean's operator type classes, and single inheritance is achieved by embedding the parent structure. See Figure \ref{lst:oop-lean} as a concrete example of a converted code. \\\\ We do not yet explore all properties of Object Oriented Programming like Polymorphism, Composition, Nested Inheritance, etc. and we leave these as part of future endeavors. See more details on implementation and extent of support of OOP in Section \ref{sec:oop}.
}

\subsubsection{Nested Definitions and Closures}

\suggest{Lean does not support nested functions similar to Python. Hence, we take a nested \code{def} and lift it to a sibling \code{private partial def} whose extra parameters are exactly the captured names (the inner reads intersected with the outer binds). This is the standard \emph{lambda lifting} transformation~\cite{johnsson1985lambdalifting}. }

\noindent
\begin{minipage}[t]{0.48\textwidth}
\begin{lstlisting}[language=python, caption={Nested Python example.}, label={lst:ll-py},  ,breaklines=true]
def make_adder(n): 
    def add(x): 
        return x + n 
    return add
\end{lstlisting}
\end{minipage}%
\hfill
\begin{minipage}[t]{0.48\textwidth}
\begin{lstlisting}[language=lean, caption={Lambda Lifted converted Lean code}, label={lst:ll-lean}, breaklines=true, mathescape=true]
private def _make_adder'add := fun (x n : Int) ↦ x +ₚ n
def make_adder := fun (n : Int) ↦ fun (x : Int) ↦ _make_adder'add x n
\end{lstlisting}
\end{minipage}



\suggest{We use sibling \code{partial def}s rather than \code{where} or \code{let rec}, since either would force the \emph{enclosing} \code{def} to be \code{partial} and cost it its unfolding and its provability(see Section \ref{app:lifting}). The lifted function needs correct set of parameters which it would have if we had used say \code{let rec}. We pass the helper exactly those variable names it reads that belong to the enclosing function, since those are the ones lifting would put out of scope. Globals and builtins need no parameter, because a sibling definition still sees them. A closure is the same lifting seen from the caller's side: a helper that escapes, by being returned or stored, becomes a partial application of the lifted definition to its captured values. One that also writes to what it captured needs a shared cell instead, since two calls to the escaped function must see each other's writes. Appendix \ref{app:lifting} gives the algorithm, the alternatives it rules out, and both closure forms. We machine-check its provable core: the free-set computation is formalised in Lean and proved scope-preserving, so lifting never puts a name out of scope that the source kept in (Appendix \ref{app:lifting}).}

\subsubsection{Multiple Returns}

\suggest{Lean requires a single return type where Python requires possibly none. We take the return type to be the join of the branch types over the lattice(Section \ref{sec:typeinfer}), so most disagreements resolve based on numeric widening principe \ref{numeric-widening}. Branches mixing \code{int} and \code{float} widen to the \code{float} type, and a branch returning \code{None} makes the result \code{Option T} rather than anything dynamic. Genuine lattice conflicts such as a \code{str} on one path and an \code{int} on another, reaches the top of the lattice, and the result is then boxed as \code{PyAny} and dispatched on its runtime tag (Section \ref{sec:pyany}). Another route is to return a proper union type, or to make the return type itself computed, giving the function a type-valued signature so that each call's return type follows from its arguments. Both are future work.}

\subsubsection{Reference Semantics and the Heap}\label{sec:heap}
Python passes mutable values by reference. Class instances, lists, dictionaries and sets are handles to
a shared object, so \code{b = a} creates a second name for the same thing and a mutation through one is
visible through the other; only immutable values --- numbers, strings, tuples --- behave as though they
were copied. The translation of Section~\ref{sec:oop} takes the second reading for everything: an
object is a Lean structure, and a mutating method returns an updated copy that is rebound at the call
site. This is faithful whenever an object has a single handle, which is the common case in the code we
target, and it keeps the program in the pure fragment where Lean's equational reasoning applies
directly. It is simply wrong once the program aliases.

We therefore provide reference semantics as a second, opt-in translation tier (\code{--heap}), rather
than changing the default, so that a program pays for the heap only when it needs one. A heap is a
partial map from addresses to values, and a translated program that uses it runs in \code{HeapM}, a
monad carrying that heap alongside Python's exceptions. Every value Python passes by reference is
allocated in the heap and represented by a typed reference \code{Ref α}; attribute access and attribute
assignment become reads and writes through that reference; and aliasing works out of the box, because
\code{b = a} now copies an address. A whole-program pass collects the types that can be stored into a
single value universe, and an interprocedural effect analysis decides which functions need the monad
at all, so pure helpers keep their pure translation. Pointer notation keeps the emitted code close to
its source.

\noindent
\begin{minipage}[t]{0.42\textwidth}
\begin{lstlisting}[language=python,caption={Python source.},label={lst:heap-py},breaklines=true]
def move(self, dx, dy):
    self.x = self.x + dx
    self.y = self.y + dy
\end{lstlisting}
\end{minipage}%
\hfill
\begin{minipage}[t]{0.54\textwidth}
\begin{lstlisting}[language=lean,caption={Generated \code{--heap} translation.},label={lst:heap-lean},breaklines=true,mathescape=true]
def Point.move (self : Ref Point) (dx dy : Int) :=
  ((do self $\rightsquigarrow$ x $\leftsquigarrow$ ($\leftarrow$ (self $\rightsquigarrow$ x)) +$_p$ dx
       self $\rightsquigarrow$ y $\leftsquigarrow$ ($\leftarrow$ (self $\rightsquigarrow$ y)) +$_p$ dy)
   : HeapM Val Unit)
\end{lstlisting}
\end{minipage}

Specifications for heap programs are separation-logic Hoare triples over the points-to assertion
($\mapsto$) and the separating conjunction ($\ast$), and Lean's \code{vcgen} discharges them
with automatic framing, so an object the callee does not touch is carried across the call with no manual reasoning.

\todo{Set up appendix with more examples and explanations}

\subsubsection{Libraries}

\suggest{Every programming language including Python heavily depends on its libraries for complex tasks, and a transpiler that handles only the core language handles very little useful code. To keep PastaLean computable, verifiable and correct, we do not use instead FFI constructs to connect Python functions within Lean since it will compromise on verifiability. We take an approach where we manually rewrite python library functions(mostly common ones) into Lean using it's standard libraries and Mathlib. It not only gives us power to verify those Libraries, but also ability to compute and verify complex codes calling those functions.}

\suggest{PastaLean supports multiple Libraries including \code{math}, \code{itertools}, \code{functools}, \code{typing}, \code{heapq}, \code{numpy}, \code{scipy}, etc. The transpiler architecture supports . We write a Lean file with the functions, and add a table that maps each Python name to the Lean name implementing it. The code generator reads that table when it sees a call from an imported module. So adding a new library, or a new function to an existing one, only means writing the definition and adding one entry to the table. The table also picks the right version for each twin, so \code{np.zeros} gives a $\mathbb{Q}$ matrix in the provable twin and a \code{Float} one in the runnable twin (Appendix \ref{app:libraries}).}

\subsubsection{IO}\label{sec:io}

Lean's \code{IO} is a grab-bag of functionalities. One monad covers files, processes, threads, mutable references,
randomness and arbitrary foreign calls, all behind an opaque token standing for the real world. That generality is the point for ordinary programming, but it makes it cumbersome for proofs. Our target programs include reading from stdin and printing to stdout, so we implemented an alternative custom monad.

The runnable twin keeps real \code{IO} and executes as usual; the provable twin instead runs in a monad combining state with Python's exceptions, where the state is a complete model of the
program's interaction with the outside world. User input is modeled an infinite stream of optional lines together
with a cursor: reading yields the cell under the cursor and advances it, and reading an absent cell throws an exception, so the read-until-end-of-input idiom terminates and \code{EOFError}
is a catchable Python exception. Output is modeled as the list of
strings printed up to any point in the execution of the program.

\subsection{Examples of Transpilation}\label{sec:transpilation}

\subsection{Dynamic Typing and Type Inference}
\label{sec:typeinfer}

Python is known for it's easy-to-use dynamic typing for users, yet a nightmare for Software Verifiers. Python supports type mutations of a variable, multiple return types of a function, it's power in Object Oriented programming, flexibility of function arguments and return types i.e. no error will be thrown if function supposed to return \code{str} returns an \code{int} type, etc. While types are dynamic, Python is strongly typed meaning it strictly enforces type rules at runtime and will raise a \code{TypeError} if incompatible types are operated on without explicit conversion. To be able to verify Python from a statically typed language like Lean, we must try to pin down the correct Type of the variable/function based it's utility if type is previously not mentioned within source code.

We propose a Type inference engine (\code{TypeInfer}) which tries to resolve Type from a given AST. The key idea is to assign to every IR node a \code{PyType} drawn from a small lattice which resolves into a Python's numeric tower into Mathlib types, and falls back to a dynamic type only when it must.

\subsubsection{The PyType Lattice}
\code{PyType} (Listing~\ref{lst:pytype_lattice}) is a finite inductive type, include a bottom \code{unknown} (``nothing known yet''), a top \code{Any} (which we call
\code{PyAny}), scalars \code{int}/\code{bool}/\code{str}/\code{float}/\code{none}, covariant
containers \code{list}/\code{set}/\code{tuple}/\code{dict}, a nullable \code{opt}, user classes
\code{cls}, and function types \code{fn}. The order has $\bot$ at \code{unknown} and
$\top$ at \code{any}, with the crucial refinement that the numeric scalars are \emph{not} flat:
they form Python's tower.

\begin{lstlisting}[caption={The \texttt{PyType} lattice.},label={lst:pytype_lattice}, mathescape=true]
inductive PyType where
  | unknown              -- $\bot$ : nothing known yet
  | any                  -- $\top$ : conflicting types, boxed as PyAny
  | int | bool | str | float | none
  | list  (elem : PyType) | set (elem : PyType)
  | tuple (elems : List PyType) | dict (key val : PyType)
  | opt   (inner : PyType)          -- Optional[T]
  | cls   (name : String)           -- a user class, e.g. TreeNode
  | fn    (args : List PyType) (ret : PyType)
\end{lstlisting}

These constructors capture the majority of types encountered in Python, ranging from primitive types to user-defined class types. Python's standard \code{typing} library provides generic type annotations such as \code{List} and \code{Dict}. However, these are normalized into the corresponding \code{PyType} constructors (e.g., \code{Dict} is represented as \code{.dict}). Each \code{PyType} constructor is internally associated with an appropriate Lean type, and operations over \code{PyType} are correspondingly mapped to their counterparts in Lean's type system.

\subsubsection{The numeric widening tower}\label{numeric-widening}
Python promotes silently: \code{True + 1} is \code{2}, \code{3 / 2} is \code{1.5}, and a list
written with both \code{5} and \code{float('inf')} is a \code{list[float]}. A naive
least-upper-bound over a flat scalar set would send \code{join(int, float)} to \code{any},
boxing the whole container and destroying provability. Instead we use the join operation which \emph{widens} along the tower. See \ref{lst:join} for the Join function used.

\[
\mathtt{bool} \;<\; \mathtt{int} \;<\; \{\mathtt{float},\ \mathbb{Q}\} \;<\; \mathbb{R},
\]

Hence \code{join(int, bool) = int} and \code{join(int, float) = float}
(Listing~\ref{lst:join}). At code-generation time the ``smaller'' operand is coerced up the
tower (\code{Int}$\to\mathbb{Q}$, and we insert the missing \code{bool} coercion rungs), rather
than the container's element type collapsing to dynamic. Containers join covariantly; tuples and
function types join pointwise when arities match and drop to \code{any} otherwise.

Type inference for a variable is performed as a monotone data-flow analysis, which can be viewed as an instance of abstract interpretation~\cite{cousot1977abstract}. Each variable initially has type \code{unknown}, and its type is gradually refined using three sources of information: (1) explicit annotations, (2) forward propagation from literals, operators, assignments, and built-ins, and (3) back-inference from how a value is used.

For example, \code{p.split()} infers \code{p : str}, while \code{x + 4} infers \code{x : int}. Similarly, \code{d[k]} provides information about the key and value types of \code{d}, and arithmetic propagates types through expressions:
\begin{lstlisting}
def f(p):
x = p.split() # p : str, x : List[str]
return x[0] # return : str
\end{lstlisting}

The analysis repeatedly propagates these constraints until no type changes. Since types only become more precise through \code{join}, this process reaches a fixed type which we call a  fixpoint. Function return types are also propagated across the module, allowing calls with known argument types to be specialized. Any type that remains \code{unknown} is finally represented as \code{PyAny}.


\subsubsection{SSA versioning across type mutations}
\label{sec:ssa}
Python rebinds freely, and a rebinding can change a variable's type:
\code{x = "5"; x = int(x)} turns a \code{str} into an \code{int}. A single Lean binder cannot
hold both. We therefore run an SSA pass that detects a
\emph{type-changing} reassignment and versions the variable (\code{x}, \code{x'v1},~\dots), with
$\phi$-joins where branches carrying different versions merge~\cite{cytron1991ssa}. Crucially,
\emph{only type changes} trigger a new version with a same-type rebinding stays a Lean \code{let mut} update, preserving the imperative shape and widening within the tower (e.g.\ \code{None} followed by an \code{int}, i.e.\ acquiring \code{Optional}) is treated as \emph{widening}, not mutation, so it does not spuriously fork the variable. Just like Standard SSA, this new variable is passed referenced instead of the previous old copies.

\subsubsection{The dynamic fallback: \texttt{PyAny}} \label{sec:pyany}

When we genuinely cannot decide a type, the join reaches \code{.any} the value is
boxed as \code{PyAny}, a tagged union of the runtime types. Following gradual
typing~\cite{siek2006gradual}, \code{PyAny} is consistent with the rest of the program: it
\emph{delegates} (rather than reimplements) the container, indexing, and truthiness protocols, so a boxed value still supports \code{x[i]}, \code{len(x)}, iteration, and \code{if x:}. However, all operations with \code{PyAny} is not possible.... \todo{complete this}


\subsection{Verification}

Whenever possible, we emit a plain Lean definition rather than a monadic program. A function whose
mutation is local can be expressed with recursion and rebinding, and is then verified like
any other Lean definition, using \code{simp}, \code{grind}, induction and Mathlib theorems. This applies to a
minority of functions, but it is the cheapest case, so we attempt it first. 

When mutation cannot be
eliminated we keep the monad as small as the program needs. If the only effect is mutability, the
translation runs in \code{Id}, whose \code{do} blocks still reduce, so a proof can recover the
functional reading. Whenever the code requires features like state, exceptions, input and output, or heap references we move the program
into the custom monads of Sections~\ref{sec:loops}, \ref{sec:io} and~\ref{sec:heap}.

\subsubsection{Contracts and Monads}

For monadic programs we use the \code{mvcgen} tactic and its successor \code{vcgen}. Given a Hoare triple goal
$\lbrace\!\lbrace P \rbrace\!\rbrace\ \mathit{prog}\ \lbrace\!\lbrace Q \rbrace\!\rbrace$, they reduce
it to a set of \emph{verification conditions}. Users can register specifications for functions so that these tactics don't unfold the definition but rather only replace calls to the function with the pre- and post-conditions from the specification. Control flow is handled through the tactics, leaving the user to provide loop invariants, termination measures, etc wherever they are needed.

The tactics require the program's monad to carry a weakest-precondition interpretation, in the form of
a \code{WPMonad} instance. Lean supplies these for \code{Id}, \code{StateT}, \code{ReaderT},
\code{ExceptT} and \code{EStateM}, but the user can also provide custom instances for other monads. \code{vcgen} additionally implements the frame rule, parameterised by a frame operator that each monad
registers along with a frame inference procedure. This is what makes separation logic usable for
our heap-manipulating programs (Section~\ref{sec:heap}).

\subsubsection{Proof Automation}\label{sec:proofautomation}

\suggest{custom tactics, AI ...}

\subsection{Limitations}

\begin{itemize}
    \item Extremely dynamic things like modifying classes at runtime
    \item Inheritance
    \item \todo{Anything else?}
\end{itemize}

\todo{Why are these OK? They're sort of marginal use cases that don't come up in the sort of scripting we're aiming this at.}

\section{Results}\label{results}

\suggest{We evaluate PastaLean on three axes: how much Python it can translate into Lean that elaborates, whether the translated program computes the same answers as the original, and how well the type inference engine recovers types. Table \ref{tab:codegen} reports the first two, Table \ref{tab:typeinfer} the third. Appendix \ref{app:methodology} describes how each benchmark was run and measured.}

\begin{table}[htbp]
\centering
\begin{tabular}{@{}lrrrr@{}}
\toprule
Dataset & Problems & Elaborated (\%) & Tests & Passed (\%) \\
\midrule
HumanEval+     & ?? & ?? & ?? & ?? \\
LeetCode       & ?? & ?? & ?? & ?? \\
LiveCodeBench  & ?? & ?? & ?? & ?? \\
\bottomrule
\end{tabular}
\caption{Code generation and execution. \emph{Elaborated} counts programs whose generated Lean type-checks; \emph{Passed} counts the dataset's own tests that the Lean twin answers identically to CPython.}
\label{tab:codegen}
\end{table}

\begin{table}[htbp]
\centering
\footnotesize
\begin{minipage}[t]{0.46\linewidth}
\centering
\textbf{Micro-benchmark}\\[2pt]
\begin{tabular}{@{}lrrr@{}}
\toprule
Category & Facts & Matched & \% \\
\midrule
return    & 230 & 169 & 73.5 \\
parameter & 95  & 62  & 65.3 \\
variable  & 525 & 314 & 59.8 \\
\textbf{all} & \textbf{850} & \textbf{545} & \textbf{64.1} \\
\bottomrule
\end{tabular}
\end{minipage}\hfill
\begin{minipage}[t]{0.46\linewidth}
\centering
\textbf{Autogen}\\[2pt]
\begin{tabular}{@{}lrrr@{}}
\toprule
Category & Facts & Matched & \% \\
\midrule
return    & 17998 & 16179 & 89.9 \\
parameter & 896   & 390   & 43.5 \\
variable  & 58374 & 39473 & 67.6 \\
\textbf{all} & \textbf{77268} & \textbf{56042} & \textbf{72.5} \\
\bottomrule
\end{tabular}
\end{minipage}
\caption{PastaLean's \code{TypeInfer} engine on the two TypeEvalPy~\cite{typeevalpy} sets (micro-benchmark left, autogen right), under the benchmark's exact-match metric. \emph{Facts} is the number of ground-truth type annotations; \emph{Matched} is those where the inferred type equals the expected one. The three categories are the benchmark's own: function return, function parameter, and local variable types.}
\label{tab:typeinfer}
\end{table}

\begin{table}[htbp]
\centering
\footnotesize
\setlength{\tabcolsep}{4pt}
\begin{tabular}{@{}lrrrr@{}}
\toprule
Static tool & Return & Param & Var & Total \\
\midrule
\textbf{PastaLean \code{TypeInfer}} & \textbf{16179} & \textbf{390} & \textbf{39473} & \textbf{56042} \\
HeaderGen        & 14086 & 346 & 36370 & 50802 \\
Jedi             & 13160 & 0   & 15403 & 28563 \\
Scalpel          & 15383 & 171 & 18    & 15572 \\
Type4Py          & 3143  & 38  & 2243  & 5424 \\
\bottomrule
\end{tabular}
\caption{Static (non-LLM) type-inference tools on the TypeEvalPy autogen set~\cite{typeevalpy} (exact matches out of 77{,}268 facts; the four baselines are the static tools of the published leaderboard). Among the leaderboard's static tools PastaLean's \code{TypeInfer} is the \textbf{strongest} --- ahead of HeaderGen, the previous static leader, by 5{,}240 facts (56042 vs.\ 50802) --- and it places 12th of 26 overall, above several instruction-tuned LLMs, while emitting a single deterministic type per fact with no model and no sampling. (Modern production type checkers, which post-date this 2024 leaderboard, are compared separately in Table~\ref{tab:checkers}; one of them, pyrefly, edges ahead on autogen through local-variable coverage.) Function-return accuracy (16179) is competitive with the strongest LLMs; the remaining gap is almost entirely local-variable coverage, and function-parameter inference is the hardest category for every tool. Per-category coverage for both TypeEvalPy sets is in Table~\ref{tab:typeinfer}.}
\label{tab:leaderboard}
\end{table}

\begin{table}[htbp]
\centering
\footnotesize
\setlength{\tabcolsep}{5pt}
\begin{tabular}{@{}lrrrrr@{}}
\toprule
Tool & Return & Param & Variable & Total & Time \\
\midrule
pyrefly            & 14246 & 54  & \textbf{46449} & \textbf{60749}\ (78.6\%) & 2.6s \\
\textbf{PastaLean \code{TypeInfer}} & \textbf{16179} & \textbf{390} & 39473 & 56042\ (72.5\%) & 9.8s \\
pyright            & 14110 & 28  & 37855 & 51993\ (67.3\%) & 11.3s \\
zuban              & 13163 & 3   & 38179 & 51345\ (66.5\%) & 1.4s \\
mypy$^\ddagger$    & 0     & 0   & 19097 & 19097\ (24.7\%) & 7.9s \\
ty$^\ddagger$      & 0     & 29  & 18942 & 18971\ (24.6\%) & 1.8s \\
\bottomrule
\end{tabular}
\caption{PastaLean vs.\ modern production type checkers on the TypeEvalPy \emph{autogen} set (77{,}268 facts; per-category denominators 17{,}998 / 896 / 58{,}374). Each checker is scored by injecting \code{reveal\_type} and matching its output name-keyed against the ground truth at root-type granularity (e.g.\ \code{list[int]}$\to$\code{list}). PastaLean leads on function \emph{return} (89.9\%) and \emph{parameter} (43.5\%) types by wide margins --- the semantically load-bearing categories --- while pyrefly's much higher local-variable coverage puts it ahead on the (variable-dominated) total; the other checkers trail. $^\ddagger$mypy and ty follow the gradual guarantee, declining to infer unannotated code by design, so they report no unannotated returns/parameters. \emph{Time} is wall-clock over the whole benchmark; PastaLean's is warm inference through its already-booted compiled backend (Mathlib load excluded), the others include process startup, so timing is indicative. On the smaller micro-benchmark (850 facts) PastaLean leads \emph{overall}: 592 (0.2s) vs.\ pyrefly 411 (0.2s), pyright 407 (1.1s), zuban 345 (0.1s), ty 214 (0.1s), mypy 205 (0.8s).}
\label{tab:checkers}
\end{table}

\clearpage

\section{Conclusions}

\small
\bibliographystyle{plain}
\bibliography{references}


\newpage

\appendix

\section{Appendices}

\subsection{Naming Conventions used in PastaLean}\label{naming-convention}

It is important to know few naming conventions used by PastaLean. To accommodate python's naming convention and to avoid possible naming conflicts, we try use Python's own naming convention which is supported by Lean. However, any extra variables/functions introduced by PastaLean use \code{'} since its usage is supported by Lean. For example, local variables introduced by us begin/end with \code{p'} , for example \code{p'_ret_1}, \code{x'v1}, etc.

For a python function \code{foo}, we name the verifiable function as \code{foo} and computable version as \code{foo'rn} suggesting it can be run. This avoids any of our naming schemes to clash with that of Python's.

While writing the theorem name's for proving a function, we use an auxiliary function proving the Hoare Triple verification as mentioned before. We name this function \code{foo_spec}. The final correctness theorem is simply \code{foo_correct}.

\subsection{Fully Converted Example}

\label{app:twins}
Figure~\ref{lst:app-src} is a Python program; Figure~\ref{lst:app-gen} is the \emph{verbatim}
Lean emitted by \code{pastalean translate --target command} in \emph{both}-twins mode (only the
\code{import}/\code{open} preamble is elided). Note how the two twins are identical except for the
carrier type: the provable \code{mean} ascribes \code{Rat} and divides exactly, while the runnable
\code{mean'rn} ascribes \code{Float} and inserts the \code{pyFloat} coercion so division matches
CPython's \code{/}. The pure \code{Id.run do} block preserves the imperative accumulator shape; the
entry point splits into a proof-mode \code{main'} (output accumulated in a pure model) and a real
\code{main''rn} in \code{IO}.

\begin{lstlisting}[style=py,caption={Python source.},label={lst:app-src}]
def mean(xs: List[int]) -> float:
    total = 0
    for x in xs:
        total += x
    return total / len(xs)

def main():
    print(mean([2, 4, 9]))
\end{lstlisting}

\begin{lstlisting}[style=lean,caption={Verbatim generated Lean (provable + runnable twins).},label={lst:app-gen}]
def mean := fun (xs : List Int) ↦
  (show Rat from
    Id.run (do
      let mut total : Int := (0 : Int)
      for x in (PastaLean.pyIter xs) do
        total := total +ₚ x
      let p'_ret_1 := total /ₚ PastaLean.pyLen xs
      return p'_ret_1))
attribute [simp, taste_ingr] mean

def mean'rn := fun (xs : List Int) ↦
  (show Float from
    Id.run (do
      let mut total : Int := (0 : Int)
      for x in (PastaLean.pyIter xs) do
        total := total +ₚ x
      let p'_ret_1 := PastaLean.pyFloat total /ₚ PastaLean.pyLen xs
      return p'_ret_1))

def main' :=
  ((do let _ ← PastaLean.ProofMode.pyPrintProof
         [pyPrintArg (mean [(2:Int),(4:Int),(9:Int)])])
   : PastaLean.ProofMode.PyProofM _)

def main''rn :=
  ((do let _ ← pyPrintIO [pyPrintArg (mean'rn [(2:Int),(4:Int),(9:Int)])]) : IO _)
\end{lstlisting}


\subsection{Provable and Computable Twin}
\todo{talk more about when you put two of them here, why you choose Array/Floats for compute and List/Q-R for Proving. Give a code snippet}

\subsection{Arithmetic: Correctness, Priorities, and Non-finite values}\label{app:arith}

Each Python operator is a single class with many instances, and Lean chooses among them from the types of the operands: $+_p$ on two strings finds the concatenation instance, on two lists the one that appends, and on an integer and a boolean the one that reads the boolean as $0$ or $1$. The result type is an \code{outParam}, so it is computed from the operands rather than supplied by the caller, and \code{1 + 2.5} yields a rational with no annotation anywhere in the generated code.

Some Python operators disagree with the instance Lean would otherwise pick, so the order in which instances are tried is itself part of the semantics. Python's \code{+} on lists concatenates, but Mathlib already defines \code{Add (List α)} pointwise; if that instance is reached first, \code{[1,2] + [3,4]} elaborates to \code{[4,6]} \ref{lst:arith-priority}, a wrong answer that no type error catches. Instances are therefore placed in one of three bands according to \emph{when} they should be consulted. A \code{high} priority instance is tried before Lean's generic ones, which is how we override cases like list concatenation. A \code{default_instance} is consulted only at the end of elaboration, when an operand's type is still open and Lean has to commit to one; ranking matters there too, because both the rational and the integer instance apply and the rational one would otherwise win, pinning an integer comprehension to $\mathbb{Q}$ and dragging the element type of the list along with it. A \code{low} priority instance is tried last, which is where the derived instances of the next paragraph sit, so that they fire only for a mixed pair nothing else covers. Listing \ref{lst:arith-priority} shows both hazards.

\begin{lstlisting}[style=lean,caption={Two places where instance ordering decides the answer.},label={lst:arith-priority}]

-- Reached first, it silently turns [1,2] + [3,4] into [4,6].
instance (priority := high) {α} : PyHAdd (List α) (List α) (List α) where
  hAdd := (· ++ ·)

-- In `[x + y for x in xs for y in ys]` the operator elaborates before the
-- type of `x` is known, so a default instance decides. Both apply, and
-- without the higher rank the ℚ one wins, forcing `xs : List ℚ`.
@[default_instance]       instance : PyHAdd Rat Rat Rat where hAdd a b := a + b
@[default_instance 10001] instance : PyHAdd Int Int Int where hAdd a b := a + b
\end{lstlisting}

Python mixes numeric types freely, so every operator must accept every pair of them, and writing one instance per operator per pair does not scale. The conversions involved do not depend on the operator, so we declare them once: \code{PyNumJoin α β γ} records that \code{α} and \code{β} both convert into a common \code{γ}, with one instance per pair in the tower $\mathbb{N} \subset \mathbb{Z} \subset \mathbb{Q}$ together with \code{Float}. A single generic instance per operator then reads those conversions off the class (Listing \ref{lst:pynumjoin}). Floor division and the bitwise operators sit outside this scheme, the first because Python's \code{//} fixes its own result type, the second because it must model unbounded integers.

\begin{lstlisting}[caption={The numeric tower, declared once and lifted to every operator.},label={lst:pynumjoin}, mathescape=true]

class PyNumJoin (α β : Type) (γ : outParam Type) where
  coeL : α → γ
  coeR : β → γ

instance : PyNumJoin Int Rat Rat := .mk2 (fun n => (n : Rat)) id

-- one generic instance per operator covers every pair in the tower
instance (priority := low) {α β γ} [j : PyNumJoin α β γ] [Add γ] : PyHAdd α β γ where
  hAdd a b := j.coeL a + j.coeR b
\end{lstlisting}

\subsubsection{Implementing \code{float('inf')}} \label{floatinf}

Python has \code{inf} to represent infinity, i.e. no integer or float is bigger than it. Programmers often use \code{float('inf')} to initialise a running minimum (correspondingly \code{float('-inf')} for a running maximum), in the idiom \code{best = inf} followed by \code{best = min(best, ...)}. In the default exact mode a Python float is a $\mathbb{Q}$, and $\mathbb{Q}$ is a field with no infinities, so there is no value for the literal to denote. For plain execution we represent it as a large finite sentinel of magnitude $10^{30}$, letting the type of the slot it lands in select the representation through a class \code{PyNonFinite}, so that a \code{Float} slot receives the genuine IEEE infinity while a $\mathbb{Q}$ or $\mathbb{Z}$ slot receives the sentinel. That runs the idiom correctly, but a finite value can never satisfy $\forall x,\ x \le \mathit{inf}$, so it carries none of infinity's order structure and supports no proof.

For verification we therefore give \code{inf} a genuine representation. \code{XRat}, defined as \code{WithBot (WithTop ℚ)} from \emph{Mathlib}, adjoins a greatest and a least element to the rationals. Because it is built from \code{Option} it stays computable, unlike \code{EReal}, which pulls in $\mathbb{R}$ and is noncomputable. \code{INF} and \code{NEGINF} are then $\top$ and $\bot$, finite values coerce in along $\uparrow$, and \code{min}, \code{max}, $\le$ and $+$ are inherited from Mathlib. The order laws the idiom depends on become theorems rather than assumptions (Figure \ref{lst:xrat}): \code{le_INF} says every value lies below \code{INF}, \code{min_INF} that \code{INF} is absorbed by \code{min}, and \code{minReduce_le} that the generated reduction is a genuine lower bound of the list it folds over. The last of these is exactly the postcondition the finite sentinel made unprovable.


\begin{lstlisting}[style=lean,caption={A provable \code{inf}: extended rationals and the laws a proof gets.},label={lst:xrat}]

abbrev XRat := WithBot (WithTop ℚ)   -- Option-based, hence computable
def INF    : XRat := ⊤
def NEGINF : XRat := ⊥

@[simp] theorem le_INF  (x : XRat) : x ≤ INF := le_top
@[simp] theorem min_INF (x : XRat) : min x INF = x := by simp

-- the generated `best = inf; for x in xs: best = min(best, x)` reduction
theorem minReduce_le (xs : List XRat) : ∀ x ∈ xs, minReduce xs ≤ x
\end{lstlisting}



\subsubsection{Correctness of the operators} \label{arith-correct}

Each custom instance for Python is definitionally Lean's own arithmetic, so we ship a proof of that correspondence alongside it. Figure \ref{lst:arith-proof} shows few examples of such correctness theorems. The proofs are trivial, but they are very important for correctness and they are what lets a tactic rewrite a generated operator back into ordinary arithmetic and then finish the goal with \code{ring}.

\begin{lstlisting}[style=lean,caption={Arithmetic Operation Proofs},label={lst:arith-proof}]

@[simp] theorem pyAdd_int (a b : ℤ) : a +ₚ b = a + b := rfl
@[simp] theorem pySub_int (a b : ℤ) : a -ₚ b = a - b := rfl
@[simp] theorem pyMul_rat (a b : ℚ) : a *ₚ b = a * b := rfl
@[simp] theorem pyDiv_rat (a b : ℚ) : a /ₚ b = a / b := rfl

\end{lstlisting}

\subsection{Why nested definitions are lifted rather than kept nested}\label{app:lifting}

Consider a nested helper whose recursion is not structural, the shape any grid or graph search takes (Listing \ref{lst:lift-why}). Lean offers three ways to keep such a helper nested, and each one fails for a different reason.

\begin{itemize}
\item A plain \code{let dfs := fun i j ↦ ...} is not recursive, so the \code{dfs} appearing in its own body is an unknown identifier.
\item \code{let rec} is recursive and does see the enclosing locals, but it cannot be marked \code{partial}, and the recursion descends toward a bound rather than on a structural subterm, so termination checking rejects it.
\item A \code{where} clause may be \code{partial}, but it inherits that \code{partial} from the enclosing definition. The outer function itself then becomes \code{partial} and loses both its unfolding and its \code{[simp]} tag, so nothing about it can be proved by unfolding. A \code{where} clause also sees only the definition's binders, and since we emit functions in the form \code{def f := fun n ↦ ...}, the parameters are lambda binders it cannot see at all.
\end{itemize}

Lifting confines \code{partial} to the helper. The enclosing function stays an ordinary \code{def}, keeps its \code{[simp]} tag, and remains both evaluable and provable.

\begin{lstlisting}[style=lean,caption={A non-structural nested recursion, and the lifted form.},label={lst:lift-why}]

def count_paths(grid):
    def dfs(i, j):
        if i == len(grid) - 1: return 1
        return dfs(i + 1, j) + dfs(i, j + 1)   # not structural
    return dfs(0, 0)

-- lifted: `partial` is confined to the helper, `grid` is passed as a capture
private partial def _count_paths'dfs := fun (i j : Int) (grid : List (List Int)) ↦ ...
def count_paths := fun (grid : List (List Int)) ↦ _count_paths'dfs 0 0 grid
\end{lstlisting}

\subsubsection{What Python and Lean change about lambda lifting}

Textbook lambda lifting~\cite{johnsson1985lambdalifting} is short: take the free variables of a nested function, make them extra parameters, move the function to the top level, and pass the captured values at every call. The transformation has itself been formalised in Lean 4 and proved correct~\cite{levy2026simplelambdalifting}; our concern is the converse, using it to \emph{produce} Lean whose enclosing definitions stay provable. That is written for a language with immutable bindings and nothing but lexical scope. Python breaks the first assumption, and its scope rules are subtle enough to have needed a tested formal semantics of their own~\cite{politz2013fullmonty}, while Lean adds requirements of its own, so we keep the shape of the transformation and change four things about how the free set is computed and two about what is emitted. Listing \ref{lst:lift-algo} shows the modified computation. 

We define two sets of variable names over a function \code{foo}. \code{usedVars(foo)} is every variable name that appears anywhere inside \code{foo}, whether it is read or written. \code{boundVars(foo)} is every variable name that \code{foo} introduces itself, that is, its parameters together with every variable it assigns to. A variable name that is used but not bound is one \code{foo} expects to get from somewhere outside itself, which is what makes it a candidate for capture. We write $A \setminus B$ for the names in $A$ that are not in $B$, $A \cap B$ for the names in both, and $A \cup B$ for their union. We write $A \mathbin{+\!\!+} B$ for $A$ followed by $B$ as a list rather than a set, which we need where the order matters: the extra parameters must be supplied at every call site in exactly the order they were added. Two smaller sets also appear in the listing: \code{comprehensionAndLambdaTargets(foo)} is the variable names that a comprehension or a lambda inside \code{foo} binds for itself, such as the \code{v} in \code{[v*2 for v in xs]}, and \code{nonlocals(foo)} is the names \code{foo} declares \code{nonlocal}.

\begin{lstlisting}[style=pseudo,caption={The free-variable computation, with the four Python-driven edits marked.},label={lst:lift-algo}]

free $\gets$ usedVars(inner) $\setminus$ boundVars(inner) // the textbook rule

// (1) a comprehension or lambda binds its own targets
free $\gets$ free $\setminus$ comprehensionAndLambdaTargets(inner)
// (2) a `nonlocal` name is assigned, so the rule above dropped it
free $\gets$ free $\cup$ nonlocals(inner)
// (3) keep only what the enclosing function itself binds
caps $\gets$ free $\cap$ boundVars(outer)
// (4) a capture the helper writes cannot travel as a parameter alone
mutated  $\gets$ every c in caps that inner rebinds or mutates in place
ordered  $\gets$ (caps $\setminus$ mutated) ++ mutated // passed in, in this order
threaded $\gets$ params(inner) mutated in place ++ mutated // handed back
\end{lstlisting}

Edit (3) is the one that keeps globals, builtins and imports out. In a purely lexical language every free name belongs to some enclosing scope, but a Python function also sees module globals and builtins, and those stay visible to a sibling top-level definition, so only a name the enclosing function itself binds would go out of scope and needs passing. Edits (2) and (4) are both consequences of Python allowing a nested function to write to an enclosing variable, which Johnsson's setting does not. A capture that is only read stays an ordinary parameter; one that is written is also handed back on the return so the caller can rebind it, which is how the write becomes visible outside while the helper stays a pure function of its inputs.

Lean forces two further changes. First, a non-structural recursion needs \code{partial}, and the nested forms spread it: a \code{where} clause inherits \code{partial} from the definition it sits in, so the \emph{enclosing} function would become \code{partial} and lose its provability. We therefore emit the helper as a sibling \code{private partial def}, which confines \code{partial} to it. Second, mutually recursive helpers have to be emitted inside one \code{mutual} block, so the nested definitions are grouped by which of them call each other and a group shares a single capture set. Lifting also runs outermost-first, since a helper nested several levels deep must capture the parameters its parent has just acquired, and each lifted parameter carries the type inferred for it in the enclosing scope, without which Lean's instance resolution has nothing to work from.

Shapes the transformation cannot express are refused rather than approximated: a helper that captures, or mutates a capture, and is then used as a value rather than called; a call inside a comprehension or lambda that would have to rebind threaded state; and a threading helper that can fall off the end without returning.

\subsubsection{Closures}

A closure is a nested function that keeps using variables from the function around it. If it only reads them we are already done, because those variables became parameters when we lifted it. Returning the function is then just calling the lifted version with the captured values filled in and the rest left open. That is what \code{make_adder} does in Listing \ref{lst:ll-lean}: \code{n} is filled in, and what comes back still takes \code{x}.

If the nested function \emph{changes} one of those variables, passing it in is not enough, because the outer function needs to see the new value. Python manages this because a closure there is a live object that keeps its variables alive, while Lean is functional and a function cannot hold state of its own, so the only way a change gets out is to return it. So the helper gives it back. It returns its own result together with the new value of the variable, and the caller stores it (Listings \ref{lst:clo-py} and \ref{lst:clo-lean}). The helper is still a plain function of its inputs, and from the outside the change looks like an ordinary assignment.

\noindent
\begin{minipage}[t]{0.44\textwidth}
\begin{lstlisting}[language=python, caption={A closure that writes to a capture.}, label={lst:clo-py}, breaklines=true]
def counter():
    count = 0
    def bump(k):
        nonlocal count
        count += k
        return count
    bump(5)
    bump(3)
    return count
\end{lstlisting}
\end{minipage}%
\hfill
\begin{minipage}[t]{0.52\textwidth}
\begin{lstlisting}[style=lean, caption={The capture is passed in and handed back.}, label={lst:clo-lean}, breaklines=true]
private def _counter'bump :=
  fun (k : Int) ↦ fun (count : Int) ↦
    let count := count +ₚ k
    (count, count)

def counter :=
  let count := (0 : Int)
  let r := _counter'bump 5 count
  let count := r.2
  let r := _counter'bump 3 count
  let count := r.2
  count
\end{lstlisting}
\end{minipage}

This works when the nested function is called in the same place it was defined. It does not work if the function both changes a captured variable and is returned, because whoever calls it later has no variable to store the new value into. That case needs the variable to become a shared cell, which is what reference semantics (Section \ref{sec:heap}) does.

\subsubsection{Machine-checked scope preservation}\label{app:lifting-proof}

Levy and Reeves prove the standard textbook algorithm correct~\cite{levy2026simplelambdalifting}. We modified that algorithm as mentioned previously, so their proof does not carry over to our version. We therefore machine-check the part we edited with Lean, that the free-set computation of Listing \ref{lst:lift-algo} preserves scope. Proving that the generated Lean behaves like the source Python would need formal semantics for both languages, which we do not have. We check that end to end by testing instead (Section \ref{results}). For the scope proof we model only what scoping depends on: the names \code{inner} uses, the names it binds, the names the enclosing function binds (\code{outerBound}), and the globals, builtins and imports ($\Gamma$) that a sibling definition can still see. We then define \code{free} and \code{caps} exactly as the four edits compute them.

\begin{lstlisting}[style=lean,caption={The free-set computation and the scope-preservation theorem (statements; proofs omitted). All are \code{sorry}-free.},label={lst:lift-proof}]
structure Inner (Name : Type) [DecidableEq Name] where
  used params assigns compTargets nonlocals : Finset Name

def boundVars (I : Inner Name) : Finset Name := I.params ∪ I.assigns
def free (I : Inner Name) : Finset Name := -- edits (1),(2)
  ((I.used \ I.boundVars) \ I.compTargets) ∪ I.nonlocals
def caps (I : Inner Name) (outerBound : Finset Name) : Finset Name := -- edit (3)
  I.free ∩ outerBound
def localBound (I : Inner Name) : Finset Name := -- genuinely-local names
  I.params ∪ (I.assigns \ I.nonlocals) ∪ I.compTargets

-- Main theorem: if inner was well-scoped nested, the lifted sibling is well-scoped.
theorem lift_preserves_scope (I : Inner Name) (outerBound Γ : Finset Name)
    (hwf : I.used ⊆ I.localBound ∪ outerBound ∪ Γ) :
    I.used ⊆ I.localBound ∪ I.caps outerBound ∪ Γ

-- Each edit does exactly its job:
theorem caps_subset_outer : I.caps outerBound ⊆ outerBound
theorem ambient_not_captured (hΓ : x ∈ Γ) (hout : x ∉ outerBound) : -- (3): globals stay out
    x ∉ I.caps outerBound
theorem nonlocal_captured : -- (2): written nonlocals in
    I.nonlocals ∩ outerBound ⊆ I.caps outerBound
theorem compTarget_not_captured (h : Disjoint I.compTargets I.nonlocals) : -- (1): comp targets out
    Disjoint I.compTargets (I.caps outerBound)
theorem ordered_eq_caps (hsub : mutated ⊆ I.caps outerBound) : -- (4): a scope-neutral reorder
    (I.caps outerBound \ mutated) ∪ mutated = I.caps outerBound
\end{lstlisting}

The main theorem says the following. Before lifting, \code{inner} can see its own locals, the enclosing bindings and $\Gamma$. After lifting, the sibling no longer sees \code{outerBound}, but it now has \code{caps}, and that is still enough to cover every name \code{inner} uses. The captures the algorithm adds replace exactly what the helper lost, so no name falls out of scope. The four smaller theorems check one edit each: globals and builtins are never captured (3), a \code{nonlocal} name that the enclosing function binds is captured (2), a comprehension target never is (1), and the split in edit (4) only reorders the same set, so it changes the calling convention and not what is in scope. 

\subsection{Libraries: a worked example and a mapping table}\label{app:libraries}

Each library needs two things. One is a Lean file that implements the functions, written using Lean's standard library and Mathlib. The other is a table that maps each Python name to the Lean name implementing it. Listing \ref{lst:numpymap} shows part of the \code{numpy} table and the registry that dispatches on the module name. The code generator itself knows nothing about \code{numpy}; it only looks the name up. Most members map to the same definition in both twins. Only a few need a different one, and those go in a small second table that is read first in exact mode. \code{np.zeros} is one of them: its element type cannot be taken from an argument, so the provable twin needs the $\mathbb{Q}$ version and the runnable twin the \code{Float} one.

\begin{lstlisting}[style=lean,caption={Part of the \texttt{numpy} table, the small exact-mode override, and the registry that dispatches on the module name.},label={lst:numpymap}]

-- one table per library: Python member name to the Lean definition
def pythonNumpyMemberMap? (member : String) : Option Lean.Name :=
  match member with
  | "array"     => some ``pyNumpyArray
  | "dot"       => some ``pyNumpyDot
  | "mean"      => some ``pyNumpyMean
  | "transpose" => some ``pyNumpyTranspose
  | "linspace"  => some ``pyNumpyLinspace
  | "zeros"     => some ``pyNumpyZerosFloat
  | _ => none

-- a few members need a different definition in the provable twin,
-- so this small table is consulted first in exact mode
def pythonNumpyMemberMapExact? (member : String) : Option Lean.Name :=
  match member with
  | "zeros" => some ``pyNumpyZerosRat
  | "ones"  => some ``pyNumpyOnesRat
  | "eye" | "identity" => some ``pyNumpyEyeRat
  | _ => none

-- one registry dispatches on the module the call came from
def pythonLibraryMap? (moduleName member : String) : Option Lean.Name :=
  match moduleName with
  | "math"  => math.pythonMathMemberMap? member
  | "numpy" => numpy.pythonNumpyMemberMap? member
  | "scipy" => scipy.pythonScipyMemberMap? member
  | _ => none
\end{lstlisting}

Listings \ref{lst:lib-py} and \ref{lst:lib-lean} show a \code{scipy} call before and after translation. \code{tmean} becomes \code{pyScipyTmean}. The arithmetic around it is translated as usual, so the library call and the code calling it are both provable.

\noindent
\begin{minipage}[t]{0.42\textwidth}
\begin{lstlisting}[language=python,caption={Python using \texttt{scipy}.},label={lst:lib-py},breaklines=true]
from scipy.stats import tmean

def variance(xs: List[float]) -> float:
    m = tmean(xs)
    total = 0.0
    for x in xs:
        total = total + (x - m) * (x - m)
    return total / len(xs)
\end{lstlisting}
\end{minipage}%
\hfill
\begin{minipage}[t]{0.54\textwidth}
\begin{lstlisting}[style=lean,caption={The generated provable twin.},label={lst:lib-lean},breaklines=true]
def variance := fun (xs : List Rat) ↦
  (show Rat from
    Id.run
      (do
        let mut m := scipy.pyScipyTmean xs
        let mut total := (0.0 : Rat)
        for x in (pyIter xs) do
          total := total +ₚ (x -ₚ m) *ₚ (x -ₚ m)
        return total /ₚ pyLen xs))
\end{lstlisting}
\end{minipage}

\subsection{Extended Semantics and Translation Rules for Python OOP} \label{oop-appendix}

This appendix provides additional details on the translation of Python's object-oriented constructs into Lean 4, expanding upon the overview in Section~\ref{sec:oop}.

Figure ~\ref{lst:oop-py} illustrates a standard Python \texttt{Point} class, while Figure ~\ref{lst:oop-lean} shows its generated Lean 4 counterpart.

\noindent
\begin{minipage}[t]{0.48\textwidth}
\begin{lstlisting}[language=python, caption={Python \texttt{Point} class.}, label={lst:oop-py}, breaklines=true]
class Point:
    def __init__(self, x, y):
        self.x = x
        self.y = y

    def norm_sq(self):
        return self.x * self.x \
             + self.y * self.y

obj = Point(1, 2)
norm = obj.norm_sq()
\end{lstlisting}
\end{minipage}%
\hfill
\begin{minipage}[t]{0.48\textwidth}
\begin{lstlisting}[language=lean, caption={Lean 4 translation.}, label={lst:oop-lean}, breaklines=true, mathescape=true]
structure Point where
  x : Int
  y : Int
  deriving Inhabited, Repr, BEq

def Point.new := fun x $\mapsto$ fun y $\mapsto$ ({ x := x, y := y } : Point)
def Point.norm_sq := fun (self : Point) $\mapsto$ self.x *$_p$ self.x +$_p$ self.y *$_p$ self.y

def obj := Point.new (1 : Int) (2 : Int)
def norm := Point.norm_sq obj
\end{lstlisting}
\end{minipage}

Decorators like \code{@static} and \code{@property} are also supported. \todo{find and mention examples.}
\todo{Talk more about limitations and what we currently support.}


\section{More on TypeInfer Engine}

In this section we talk more about the Type Inference and some intricacies to deal with while modeling Python. Note that the Type Infer engine has been refined as it was tested more and more with different and interesting pieces of code.

\subsection{Join over PyType Lattice}\label{join-pytype}

As mentioned in \ref{app:lattice}, we use the \emph{join} function to go up the numeric widening tower maintaining soundness of the type checker. These join operations are made keeping in mind behavior of Python.

\begin{lstlisting}[language=lean,caption={Widening join over the numeric tower (excerpt).},label={lst:join}]
def join : PyType → PyType → PyType
  | .unknown, t | t, .unknown => t  -- ⊥ is identity
  | .any, _ | _, .any => .any  -- ⊤ absorbs
  | .int, .bool | .bool, .int => .int  -- bool < int
  | .float, .int | .int, .float
  | .float, .bool | .bool, .float => .float  -- int,bool < float
  | .list a, .list b => .list (join a b)  -- covariant
  | .opt a, b | b, .opt a => match join a b with
                    | .any => .any | j => .opt j  -- Optional[Any] ≡ Any
  | .none, t | t, .none => .opt t
  | a, b => if a.beq b then a else .any
\end{lstlisting}

A subtle point that took real care is that \code{opt}/\code{none} must \emph{absorb} into \code{any}
(``a nullable dynamic value is just a dynamic value''). Without the rule
\code{join (.opt a) b} that collapses \code{Optional[Any]} to \code{Any}, \code{join} is
non-associative \code{join(join(None,int),str)} and \code{join(None,join(int,str))} disagree
(\code{Option PyAny} on one grouping, \code{PyAny} on the other) which makes the predicted Type non-confluent. Associativity of \code{join} is exactly what a lattice fixpoint needs, so we
restored it by construction.

\subsection{The lattice laws are machine-checked}
\label{sec:latticeproof}
The soundness of inference rests on properties of \code{join} and \code{consistent}: that
\code{join} is a genuine bounded join-semilattice (so the fixpoint terminates and is
order-independent) and that \code{consistent} is a gradual-typing consistency relation. We do not
merely assert these---we \emph{prove} them in Lean, with no \code{sorry}, over a finite model that keeps one representative per distinct
join behavior (so the equations reduce by computation). We establish \code{join\_comm},
\code{join\_assoc}, \code{join\_idem}, and the bounds \code{join\_unknown} ($\bot$) /
\code{join\_any} ($\top$); that the induced order ($a\le b$ iff \code{join a b} $=b$) is a partial
order with \code{join} as least upper bound and \emph{monotone} (the property that makes the
reflow converge); and, for the dynamic type, \code{consistent\_refl}, \code{consistent\_symm}, and
crucially \code{consistent\_not\_trans}---the \emph{failure} of transitivity
(\code{int}$\sim$\code{any}$\sim$\code{str} yet \code{int}$\not\sim$\code{str}), which is exactly
what distinguishes gradual consistency~\cite{siek2006gradual} from subtyping. The production
functions are then pinned to the model by \code{\#guard}/\code{native\_decide}. To our knowledge
this is an unusual property for a transpiler: its type system's meta-theory is verified in the
same assistant it targets.

The metatheory of \S\ref{sec:latticeproof} is established over a
finite model \code{Ty} (one representative per distinct join behavior), with no \code{sorry}:

\begin{itemize}\itemsep1pt
\item \textbf{Bounded join-semilattice.} \code{join\_comm}, \code{join\_assoc}, \code{join\_idem};
bounds \code{join\_unknown} ($\bot$ is identity) and \code{join\_any} ($\top$ absorbs).
\item \textbf{Induced order.} ($a \le b$ iff \code{join a b} $= b$) is a partial order
(\code{le\_refl}, \code{le\_trans}, \code{le\_antisymm}); \code{le\_join\_left}/\code{le\_join\_right}
give monotonicity---the property that forces the reflow fixpoint to converge---and \code{join\_lub}
proves \code{join} is the least upper bound.
\item \textbf{Gradual consistency.} \code{consistent\_refl}, \code{consistent\_symm},
\code{consistent\_unknown} (the gradual guarantee), and \code{consistent\_not\_trans}: a machine-checked
witness that consistency is \emph{not} transitive, the property separating gradual typing from
subtyping.
\item \textbf{Coercions total.} \code{reconcile\_refl} and \code{reconcile\_total}: every implicit
coercion the code generator may insert is one of the finite intended actions
(\code{boolToInt}, \code{intToFloat}, \code{unwrapOpt}, \code{box}), never ``stuck.''
\end{itemize}
The production \code{join}/\code{consistent}/\code{reconcile} used by the transpiler are then pinned
to this model by \code{\#guard}/\code{native\_decide} in \code{PALC/TypeInfer/TestLattice.lean}.
% TODO(author): paste the full `join`/`consistent`/`reconcile` and the associativity proofs verbatim.


\subsection{The dynamic fallback: \code{PyAny}}\label{app:pyany}

When inference cannot give a value a single Lean type, it is boxed as \code{PyAny}, a tagged union of the runtime types. Boxing happens at the boundary through a \code{CoeTail} instance, so no explicit wrapper appears in the generated code: a function returning an \code{int} on one path and a \code{str} on another coerces on both. The same fallback absorbs a variable rebound to a different type midway through a function, and a \code{try}/\code{except} whose branches return different types. Listing \ref{lst:poly-any} shows one \code{add} applied at two types.

\begin{lstlisting}[style=lean,caption={One \code{add} used at two types: the parameters box to \code{PyAny}, and a property of the boxed function is still proved.},label={lst:poly-any}]

# Python: the call sites disagree, so the join of int and str is the top
def add(a, b):    return a + b
print(add(3, 4)); print(add("x", "y"))

-- generated Lean: both parameters are boxed
def add := fun (a : PyAny) ↦ fun (b : PyAny) ↦ (show PyAny from a +ₚ b)

-- and associativity of the boxed operator is still provable
example (a b : PyAny) : (a +ₚ b) +ₚ b = a +ₚ (b +ₚ b) := by
  pyany_cases a b <;> grind
\end{lstlisting}

Boxing costs provability, since \code{PyAny} is not a commutative ring and \code{ring} stalls on a boxed goal. The tactic \code{pyany_cases} recovers much of it by splitting the goal into one goal per runtime type the value could hold, discharging the tag-mismatch leaves automatically. The recovery is partial by nature: $+_p$ is not commutative over \code{PyAny}, since it concatenates on strings and lists, and the tactic correctly leaves that goal open.


\section{Evaluation Methodology}\label{app:methodology}

\todo{Describe the harness, the machine, timeouts, and how a run is reproduced.}

\subsection{LeetCode and LiveCodeBench}\label{app:eval-cp}

\todo{How problems and solutions are selected, how tests are executed, what counts as a pass.}

\suggest{\paragraph{LeetCode}: The Leetcode dataset contains \textbf{2{,}589} LeetCode problems. More than \textbf{92\%} ($\approx$2{,}389) transpile and compile, and of those, the runnable twin reproduces CPython's answer on \textbf{$>$95\%} of test cases. The fail tests include primarily Timeout's from Lean's compiled code due to internal time complexity of implementation of Data Structures as highlighted in Section TODO.}

\suggest{\paragraph{LiveCodeBench}: The LiveCodeBench contains various sub datasets for execution, code generation, etc. However for our use case, we use the \code{execution-v2} latest dataset for Code Generation from Python GroundTruths containing \textbf{476} problems. More than \textbf{96.6\%} ($\approx$460) transpile and compile, and of those, the runnable twin reproduces CPython's answer on \textbf{$>$99.2\%} of test cases. }

\subsection{Type Inference}\label{app:eval-type}

\paragraph{TypeEvalPy.} TypeEvalPy~\cite{typeevalpy} is a type-inference \emph{recall} benchmark (ICSE 2024): every ground-truth type annotation is a \emph{fact}, keyed by source location and name, and a tool is scored by how many facts it recovers by exact string match after both sides are normalised to a common vocabulary (\code{int}, \code{str}, \code{list}, \code{Callable}, class names, \code{None}, \ldots). Facts fall into three categories --- function return, function parameter, and local variable types. It ships two sets: a hand-written \emph{micro-benchmark} (153 snippets, 850 facts across 18 Python-feature categories) and the much larger \emph{autogen} set (5{,}453 snippets, 77{,}268 facts) generated by expanding the micro templates. To score \code{TypeInfer} we run the engine on each snippet, read the \code{PyType} it assigns to each annotated location, lower it to the benchmark's vocabulary, and compare against the ground truth; \emph{covered} counts facts we assign any type, \emph{matched} counts exact matches. We report both in Table~\ref{tab:typeinfer}, and place \code{TypeInfer} against the published leaderboard in Table~\ref{tab:leaderboard}.

On the autogen set \code{TypeInfer} matches 56{,}042 / 77{,}268 facts (72.5\%), which makes it the \textbf{best static (non-LLM) tool} on the leaderboard --- ahead of HeaderGen, the previous static leader, by 5{,}240 facts --- and 12th of 26 overall, above several instruction-tuned LLMs, while emitting a single deterministic type per fact with no model and no sampling. Function-return inference (16{,}179, 89.9\%) is competitive with the strongest LLMs; the remaining gap is dominated by local-variable coverage, and function-parameter inference (43.5\%) is the hardest category for \emph{every} tool on the board. On the smaller micro-benchmark it matches 545 / 850 (64.1\%). The published leaderboard covers only the autogen set, so Table~\ref{tab:leaderboard} is over autogen; the micro numbers are PastaLean's own run.

\subsection{HumanEval}\label{app:eval-humaneval}

\todo{Contract and non-contract variants, and the eleven-problem note.}

\suggest{\paragraph{HumanEvalPlus}: On the 164 HumanEval~\cite{chen2021codex} tasks, PastaLean transpiles and compiles \textbf{153/164} with \emph{no} human annotation and $0$ degraded (\code{pyUnsupported}) statements. Meanwhile the remaining 11 compile once given a light type hint on a single parameter (their local types are genuinely undecidable, e.g.\ a parameter used as both \code{str} and \code{list[str]}). Thus for decidable type's, our code compiles and runs \textbf{164/164}. Executing the runnable twins against the HumanEval$^+$ test suites, the outputs match CPython on \textbf{98.9\%} of test cases (4{,}773/4{,}827), with \textbf{95} solutions fully correct on every case.}

\end{document}